% abstract.tex: BORRADOR (2026-10-18), traducción de resumen.tex; revisar junto con él.
This report evaluates the viability of variational quantum transformations as an encoding mechanism for morphological features of real mammograms. On the 3,568 lesions of CBIS-DDSM, with masses and calcifications treated as independent subproblems, 67 radiomic features are extracted with PyRadiomics, 12 are selected and scaled to $[0,\pi]$. The same vector feeds five conditions: no transformation, PCA, RBF Kernel PCA, and two quantum feature maps (ZZ, and Pauli with Y and ZZ terms) followed by a fixed-parameter ansatz, simulated with Qiskit Aer. Separability is measured with kernel-target alignment, regularized geometric difference, effective dimension, the Davies-Bouldin index and the Fisher distance; classification, with the same multilayer perceptron and patient-grouped cross-validation.

The quantum representations show no greater separability in any metric. The quantum kernel is concentrated (off-diagonal median close to 0.003, against 0.61 for the RBF kernel) and the geometric difference (1.40 to 1.70) stays well below the $\sqrt{N}\approx 14$ threshold. The perceptron loses between 0.07 and 0.14 AUC-ROC on the quantum embeddings, while an SVM on the quantum kernel matches the classical one, which places the loss in the expectation-value readout rather than in the encoding. The simulator constraints are also quantified: with 1,024 shots, between 16 and 74\,\% of the kernel entries would read exactly zero.

\medskip
\noindent\textbf{Keywords:} quantum machine learning, quantum embeddings, quantum kernels, class separability, radiomics, CBIS-DDSM, breast cancer.
